\scanpage{97}
\mspage{73}
Again note that $\mathbb S_x\mathbb H_xJ$ is represented by the matrix whose elements are the symbols of theelements of the matrix corresponding to $\mathbb H_xJ$. Thus, if $\xi(x)$ is represented by $(\xi^{ij})$ and $F_x$ by $(F^{ij})$, then we obtain

\uline{Theorem 4.6.8.} If $F^{ij}=\displaystyle\sum_{n=1}^{\infty}Y_n^{ij}$, where $Y_n^{ij}$ is a spherical harmonic of degree $n$, then
\[
(\mathbb S_x\mathbb H_xJ)(u)=\xi^{ij}+\sum_{n=1}^{\infty}\gamma_nY_n^{ij}(u),
\quad u\in S_x
\]
where
\[
\gamma_n=i^n\pi^{\frac m2}\frac{\Gamma(\frac n2)}{\Gamma(\frac{m+n}{2})}.
\]

If $J\in\mathcal G^p$, we again say that $J\rho=\alpha$ is regularizable if there exists $K\in\mathcal G^p$ such that $KJ\rho=K\alpha$ is a Fredholm equation with the same solutions as $J\rho=\alpha$. Now by applying the results of systems of singular equations, see [20f] pp. 106-109, we get

\uline{Theorem 4.6.9.} If $J\in\mathcal G^p$ then the equation $J\rho=\alpha$ is regularizable if and only if
$\det[\mathbb S_x\mathbb H_xJ(u)]\ne0$ for all $u\in S_x$ and all $x\in M$.

Now let $\mathcal S_M^p$ be the algebra of bundle maps of $S(M)\longrightarrow A^p\otimes A^p$. We define a map
$\mathbb S_M:\mathcal G^p\longrightarrow\mathcal S_M^p$ by
\mspage{96}
